\documentclass[11pt,aspectratio=169]{ctexbeamer}

\usetheme{default}
\setbeamertemplate{navigation symbols}{}

\usepackage{amsmath}
\usepackage{booktabs}

\definecolor{DeepBlue}{RGB}{25,74,122}
\definecolor{SoftInk}{RGB}{53,59,66}
\definecolor{PaperBg}{RGB}{249,248,245}
\definecolor{PanelBlue}{RGB}{236,243,249}

\setbeamercolor{normal text}{fg=SoftInk,bg=PaperBg}
\setbeamercolor{structure}{fg=DeepBlue}
\setbeamercolor{frametitle}{fg=DeepBlue,bg=PaperBg}
\setbeamercolor{title}{fg=DeepBlue}
\setbeamercolor{block title}{fg=DeepBlue,bg=PanelBlue}
\setbeamercolor{block body}{fg=SoftInk,bg=white}

\newcommand{\Core}[1]{
  \begin{block}{核心结论}
    #1
  \end{block}
}

\newcommand{\LectureMeta}[2]{
  \title{大模型推理基础设施}
  \subtitle{#1}
  \author{张书豪}
  \institute{华中科技大学计算机学院}
  \date{#2}
}


\LectureMeta{第 12 讲 开源系统实践}{研究生课程课件}

\begin{document}

\begin{frame}[plain]
  \titlepage
\end{frame}

\begin{frame}{教学目标}
  \Core{本讲将前述概念落实到真实开源 runtime、benchmark 与复现实验流程之中。}
  \begin{itemize}
    \item 学会进入真实代码库。
    \item 学会定位请求路径、KV 路径和指标输出。
    \item 学会记录最小复现实验。
  \end{itemize}
\end{frame}

\begin{frame}{进入代码库的最小路径}
  \Core{进入系统的关键不在于覆盖全部代码，而在于准确定位入口路径。}
  \begin{itemize}
    \item 请求入口在哪里。
    \item 调度循环在哪里。
    \item KV 分配与回收在哪里。
    \item benchmark 和结果输出在哪里。
  \end{itemize}
\end{frame}

\begin{frame}{代码导读的方法}
  \Core{代码导读的关键在于将核心路径串联起来，而非简单扩大阅读文件数量。}
  \begin{itemize}
    \item 先追请求入口。
    \item 再追调度循环和 batch 组织。
    \item 最后追 KV 与 metrics 输出点。
  \end{itemize}
\end{frame}

\begin{frame}{最小代码导读任务}
  \Core{一次合格的代码导读，应至少能把请求、状态和指标三条路径串起来。}
  \begin{itemize}
    \item 从 API 到 scheduler
    \item 从 scheduler 到 model runner
    \item 从 model runner 到 KV manager 和 metrics
  \end{itemize}
\end{frame}

\begin{frame}{最小复现实验任务}
  \Core{复现实验的关键不在于跑出数字，而在于能解释数字是如何产生的。}
  \begin{itemize}
    \item 记录模型、输入、参数和硬件环境。
    \item 记录 workload 构造方式和 warmup 过程。
    \item 保留日志、脚本和结果表格。
  \end{itemize}
\end{frame}

\begin{frame}{源码穿插：nano-vllm-hust 作为课程实践载体}
  \Core{\texttt{\detokenize{nano-vllm-hust}} 已经按教学目的组织为“参考实现 + 学生练习”双层结构，适合作为本课程的实践载体。}
  \begin{itemize}
    \item \texttt{\detokenize{nanovllm/}} 保留可运行参考实现，适合课堂演示和源码导读。
    \item \texttt{\detokenize{student_exercises/}} 提供带 TODO 的练习版本，适合课程实验。
    \item \texttt{\detokenize{TEACHING.md}} 给出了 tutorial 与源码文件之间的映射关系。
  \end{itemize}
\end{frame}

\begin{frame}{源码穿插：建议的课堂实践顺序}
  \Core{基于该代码库，可以把课堂实践组织为由浅入深的五步式路径。}
  \begin{itemize}
    \item 第一步：从 \texttt{\detokenize{llm_engine.py}} 跟踪请求进入与单步推进。
    \item 第二步：在 \texttt{\detokenize{scheduler.py}} 中分析调度与回收逻辑。
    \item 第三步：在 \texttt{\detokenize{block_manager.py}} 中理解 KV 分配、复用与哈希。
    \item 第四步：结合 \texttt{\detokenize{bench.py}} 讨论如何形成可复现实验记录。
  \end{itemize}
\end{frame}

\begin{frame}{源码穿插：作业包如何组织}
  \Core{`nano-vllm-hust` 已提供面向 tutorial 的作业打包脚本，可将课堂实践直接转化为标准化作业目录。}
  \begin{itemize}
    \item \texttt{\detokenize{scripts/teaching/prepare_assignment.py}} 用于生成每次 tutorial 的作业包。
    \item 生成结果位于 \texttt{\detokenize{assignments/tutorial_X_spring-2026/}}。
    \item 每个作业包同时包含说明文件、starter 文件与提交模板。
  \end{itemize}
\end{frame}

\begin{frame}{复现实验中的常见遗漏}
  \Core{复现失败往往来自配置、默认参数与脚本细节记录不完整。}
  \begin{itemize}
    \item 模型版本和推理精度。
    \item warmup、batch 限制和长度分布。
    \item benchmark 的默认请求生成逻辑。
  \end{itemize}
\end{frame}

\begin{frame}{实践环节的重要性}
  \Core{进入真实系统后，才能观察到研究问题常常位于机制设计与工程边界的连接处。}
  \begin{itemize}
    \item 论文中的一个框图，在系统里可能跨越多个模块。
    \item 实验结果往往依赖脚本、默认参数和日志解释。
    \item 这也是课程项目最有价值的训练部分。
  \end{itemize}
\end{frame}

\begin{frame}{例子：进入一个开源 runtime 的正确姿势是什么}
  \Core{开源系统实践课最怕学生一上来就“全仓库全局搜索”，最后只积累了碎片印象。}
  \begin{block}{错误路径}
    先把目录全翻一遍，却没有确定任何主调用链和控制点。
  \end{block}
  \pause
  \begin{block}{正确路径}
    先找入口，再找调度循环，再找状态对象，最后找指标与脚本出口。
  \end{block}
  \pause
  \begin{block}{课程结论}
    代码实践的第一原则不是多看，而是先抓住最短的真实路径。
  \end{block}
\end{frame}

\begin{frame}{例子：为什么 benchmark 脚本也必须纳入源码导读}
  \Core{系统知识不只在 runtime 本体里，benchmark 脚本经常在定义 workload、指标和报告方式。}
  \begin{block}{runtime 负责执行}
    它告诉你请求如何被调度、状态如何被组织、结果如何被产生。
  \end{block}
  \pause
  \begin{block}{benchmark 负责解释}
    它告诉你输入是什么、统计口径是什么、图表为什么长成这样。
  \end{block}
  \pause
  \begin{block}{教学意义}
    不读 benchmark，就很容易误把“脚本默认值”当成“系统本身结论”。
  \end{block}
\end{frame}

\begin{frame}{代码例子：最小代码导读清单}
  \Core{开源系统实践课适合给学生一份很短但很硬的导读清单。}
  \begin{block}{导读顺序}
    {\ttfamily\small
    \detokenize{1. LLMEngine.generate()}\par
    \detokenize{2. Scheduler.schedule()}\par
    \detokenize{3. BlockManager.allocate()}\par
    \detokenize{4. bench.py / result writer}\par
    }
  \end{block}
  \pause
  \begin{block}{课堂应追问}
    为什么这四步足以构成一个最小系统理解闭环？
  \end{block}
  \pause
  \begin{block}{结论}
    好的课堂实践不会让学生读完所有代码，而是让他们先得到一条可复述的主路径。
  \end{block}
\end{frame}

\begin{frame}{代码例子：一次最小复现实验记录应该写什么}
  \Core{“我跑通了”不构成复现；复现必须留下别人可重复的上下文。}
  \begin{block}{记录模板}
    {\ttfamily\small
    \detokenize{model=..., precision=..., max_model_len=...}\par
    \detokenize{workload=..., warmup=..., num_requests=...}\par
    \detokenize{command=..., logs=..., result_table=...}\par
    }
  \end{block}
  \pause
  \begin{block}{课堂应追问}
    哪些字段缺失时，哪怕数字是对的，也不能算可复现？
  \end{block}
  \pause
  \begin{block}{课程结论}
    复现实验文档是工程记录，也是研究证据的一部分。
  \end{block}
\end{frame}

\begin{frame}{例子：为什么学生最容易在默认参数上踩坑}
  \Core{很多所谓“复现失败”并不是机制没实现，而是默认值、环境与隐藏脚本分支没有被看见。}
  \begin{block}{常见遗漏}
    温启动、随机种子、batch 上限、长度分布和精度模式。
  \end{block}
  \pause
  \begin{block}{常见误判}
    看到不同数字就直接怀疑代码逻辑，而不是先检查实验边界是否一致。
  \end{block}
  \pause
  \begin{block}{教学重点}
    进入真实开源系统后，最需要养成的是“先核对边界，再解释数字”的习惯。
  \end{block}
\end{frame}

\begin{frame}{课堂练习：给自己设计一条 30 分钟导读路线}
  \Core{实践课的目标不是一口吃掉系统，而是能在短时间内稳定定位关键对象。}
  \begin{block}{练习一}
    规定 30 分钟时间，写出你会先打开的 4 个文件和顺序。
  \end{block}
  \pause
  \begin{block}{练习二}
    说明每个文件分别是为了解决“入口、控制、状态、证据”中的哪一个问题。
  \end{block}
  \pause
  \begin{block}{练习三}
    用一句话回答：为什么 benchmark 脚本必须算作系统知识的一部分？
  \end{block}
\end{frame}


\begin{frame}{课堂讨论}
  \Core{实践环节的目标在于建立可解释、可修改、可验证的系统理解。}
  \begin{itemize}
    \item 什么是进入一个 runtime 的正确起点？
    \item 为什么 benchmark 脚本本身就是系统知识的一部分？
    \item 复现实验最容易遗漏哪些导致结论失真的细节？
  \end{itemize}
\end{frame}

\begin{frame}{本讲总结}
  \Core{通过代码与实验进入真实系统，是将课程知识转化为研究能力的关键步骤。}
  \begin{itemize}
    \item 项目后续将以这一实践基础继续推进。
    \item 下一讲将进入项目问题定义和工作坊环节。
  \end{itemize}
\end{frame}

\end{document}